\icode{E-T-SI-L1-260625-002}
将一个分子\(P\)固定在\(O\)点，另一个分子\(Q\)从图中的\(a\)点由静止释放，两分子之间的作用力与间距关系的图象如图所示，则下列说法正确的\paren[]

\begin{tikzpicture}[>=latex,font=\scriptsize]
    \draw[->,blue] (0,-1.5)--(0,2)node[left]{\(F\)};
    \draw[->,blue] (0,0)node[left]{\(O\)}--(4,0)node[below]{\(r\)};
    \draw[purple] (.3,1.8)..controls(.6,-1.5)and(1,0)..(3.6,-.02);
    \filldraw (.63,0)circle(.04);
    \draw[dashed] (1.13,0)--(1.13,-.3);
    \node at(.69,0)[above]{\(c\)};
    \filldraw (1.13,0)node[above]{\(b\)} circle(.04);
    \filldraw (2.5,0)node[above]{\(a\)}circle(0.04);
\end{tikzpicture}
\begin{choices}
    \item 分子\(Q\)由\(a\)运动到\(c\)的过程中，先加速再减速
    \item 分子\(Q\)在\(b\)点时分子势能最小
    \item 分子\(Q\)在\(c\)点时加速度大小为零
    \item 分子\(Q\)由\(a\)点释放后运动到\(c\)点左侧的过程中，加速度先增大后减小
\end{choices}
